% ECONOMICS_PROBLEM: {"id": "D63-242", "title": "EF1 existence for general mixed-item utilities", "jel_code": "D63", "source_id": "OP-0189"}
% !TeX program = xelatex
\documentclass[11pt,a4paper]{article}
\usepackage[margin=23mm]{geometry}
\usepackage{fontspec}
\setmainfont{Latin Modern Roman}
\usepackage{amsmath,amssymb,mathtools}
\usepackage{xcolor,enumitem,booktabs,longtable,array,xurl,hyperref}
\definecolor{muted}{HTML}{53636C}
\hypersetup{colorlinks=true,urlcolor=blue,linkcolor=blue}
\setlength{\parindent}{0pt}
\setlength{\parskip}{5pt}
\newcommand{\R}{\mathbb R}
\newcommand{\E}{\mathbb E}
\newcommand{\N}{\mathbb N}
\newcommand{\ind}{\mathbf 1}
\newcommand{\Var}{\operatorname{Var}}
\newcommand{\Cov}{\operatorname{Cov}}
\newcommand{\supp}{\operatorname{supp}}
\DeclareMathOperator*{\argmin}{arg\,min}
\DeclareMathOperator*{\argmax}{arg\,max}
\newcommand{\entrymeta}[3]{{\sffamily\small\color{muted}\textbf{\texttt{#1}}\quad|\quad #2\par #3\par}}
\newcommand{\Problemref}[1]{\texttt{#1}}
\newcommand{\JELref}[2]{\texttt{#2} (JEL \texttt{#1})}
\begin{document}
\section*{EF1 existence for general mixed-item utilities}
\textbf{Problem ID:} \texttt{D63-242}\quad \textbf{JEL:} \texttt{D63}\par
% BEGIN ECONOMICS STATEMENT
\label{problem:OP-0189}
% GAME_THEORY_PROBLEM: {"local_id": "GT-059", "original_number": 59, "kind": "P", "entry_kind": "statement_only", "published": false, "review_required": true, "category": "game_theory"}
\label{game:GT-059}
{\sffamily\small\color{muted}
\textbf{GT-059} | Mixed fair division | \textbf{P}: source-date general-utility question with a fixed EF1 convention.
\par}
\subsection*{Definitions and mathematical statement}
Let $N=[n]$, $n\geq3$, $M$ finite, and $u_i:2^M\to\mathbb R$ be normalized functions with $u_i(\varnothing)=0$. No additivity or monotonicity is assumed. A complete allocation partitions $M$. In the one-deletion mixed EF1 convention, for every ordered pair $i,j$ at least one of the following holds:
\[
 \begin{split}
 &u_i(A_i)\geq u_i(A_j),\quad\text{or}\\
 &\exists g\in A_j:\ u_i(A_i)\geq u_i(A_j\setminus\{g\}),\quad\text{or}\\
 &\exists g\in A_i:\ u_i(A_i\setminus\{g\})\geq u_i(A_j).
 \end{split}
\]
Determine whether every such profile admits a complete allocation satisfying these conditions. Also isolate the subquestion $u_1=\cdots=u_n$: does the conclusion hold for every identical normalized utility function?
\subsection*{Short English statement}
Can envy always be eliminated by removing at most one item from either of the two compared bundles?
\subsection*{Variants and scope boundaries}
With arbitrary nonmonotone functions, an item need not have a bundle-independent good/chore sign; the displayed inequalities define the intended notion directly. Additive mixed utilities are a known positive subcase. Removing one item from each bundle simultaneously is a different, weaker convention.
\subsection*{Source relationship and status}
The survey's Open Question 1 explicitly asks this general-utility question, including identical utilities. Definition 3.3 supplies the separate alternative deletions. The question is retained at the retrieved source version date; later status has not been exhaustively certified.
\subsection*{References}
\begingroup\footnotesize\raggedright
\textbf{[S1]} Shengxin Liu, Xinhang Lu, Mashbat Suzuki, and Toby Walsh (2024 version; first posted 2023). \emph{Mixed Fair Division: A Survey}. arXiv:2306.09564v4. \textit{Verified locator:} Definition 3.3; Section 4.1, Open Question 1. \url{https://arxiv.org/html/2306.09564v4}
\par\endgroup

\subsection*{Common conventions: Source mathematical conventions}

\textbf{Finite games.} For a finite set $A$, $\Delta(A)$ is its probability
simplex. Players choose independent mixed strategies in Nash equilibrium;
correlation is allowed only when explicitly part of the solution concept.
In a game with payoff functions $u_i$ and strategy profile $x$, an additive
$\varepsilon$ Nash equilibrium satisfies
\[
 u_i(x)\geq u_i(x_i',x_{-i})-\varepsilon
 \quad\text{for every player $i$ and every admissible deviation $x_i'$}.
\]
For numerical approximation, payoffs are normalized as locally specified.
An $\varepsilon$ payoff residual is distinct from distance at most
$\varepsilon$ to the set of exact equilibria. Point convergence, convergence
of empirical play, and convergence in distance to an equilibrium set are
also distinct.

\textbf{Computational representations.} Unless stated otherwise, a complexity
question takes rational data with binary encoding; polynomial time is measured
in total bit length. A fixed-accuracy polynomial algorithm is not necessarily
polynomial in $1/\varepsilon$, and neither is necessarily polynomial in
$\log(1/\varepsilon)$. Succinct games and value, demand, or sampling oracles
must declare their interfaces. Exact real solutions require an explicit
representation; an unspecified list of arbitrary real numbers is not a
finite computational output. A claimed algorithmic barrier is meaningful
only for its specified model and reductions.

\textbf{Dynamic games.} Histories, information, observation, and allowable
randomization are part of the model. A stationary strategy depends only on
the current state; a positional strategy in a deterministic graph game
chooses one action at each controlled vertex. Nash equilibrium at an initial
state and equilibrium after every history are different requirements.
An expectation of a pathwise $\liminf$ payoff, a $\liminf$ of expected averages,
a limiting expected average, and a uniform finite-horizon equilibrium are
different criteria. None is silently substituted for another.

\textbf{Allocation and mechanisms.} A complete allocation assigns every item
exactly once unless otherwise stated. Zero-valued goods, negative values,
capacity constraints, feasibility, lotteries, and copies of goods affect
fairness definitions and are specified locally. Dominant-strategy incentive
compatibility, Bayesian incentive compatibility, universal truthfulness,
and truthfulness in expectation impose different inequalities. Whenever
an objective ratio has zero denominator, its convention is stated or those
instances are excluded.

\textbf{Infinite-dimensional models.} Expectations, integrals, stochastic
equations, and optimization objectives require their local regularity,
integrability, filtrations, and admissible domains. A backward--forward
mean-field system is a boundary-value problem; it is not an initial-value
semiflow without an additional construction. Long-time, large-population,
and vanishing-noise limits require an order of limits and a convergence
topology. If a minimum need not be attained, the objective is its infimum
or supremum and the approximation requirement is stated separately.
% END ECONOMICS STATEMENT
\end{document}
